不同的力，做功的性质可能很不相同。一类是重力、理想弹簧弹力等力：它们所做的功只与
物体的初、末位置有关，与物体在力作用下经过的轨迹无关。另一类是滑动摩擦力等各种阻力：
它们所做的功与实际经过的路程有关。阻力会消耗机械能，工程设计中常设法减小它，但在真实
环境中又难以完全避免。当物体在无阻力的空间中运动，或者阻力小到可以忽略时，运动会表现
出一个十分清楚的共同性质。但“没有阻力”还不够：只有保守力作功、其他力不作功时，才有本节讨论的机械能守恒。

\subsection{重力作用下物体的运动}

如图~\ref{ch08:fig:gravityPathsDiagram}，地球表面附近有高度差为 $h$ 的两点 $A,B$。在两点之间设置若干形状
不同、固定不动的光滑轨道，让物体从 $A$ 点沿轨道运动。以下比较以物体确能到达 $B$、途中保持约束为前提；
若路径先升高，初动能还须足够越过最高处。轨道对物体的压力始终垂直于轨道切线，
因而不做功；重力功只取决于初、末位置的高度差：
\gravityPathsDiagram
\begin{equation}
 W_G=mgh .
\end{equation}
设所有物体在 $A$ 点的初速率均为 $v_A$，由动能定理
\begin{equation}
 mgh=\frac12mv_B^2-\frac12mv_A^2.
\end{equation}
移项得到
\begin{equation}
 \frac12mv_B^2=\frac12mv_A^2+mgh.\label{ch08:eq:gravity-kinetic}
\end{equation}
虽然物体在 $B$ 点的速度方向受到各自轨道的限制，但它们的速率和动能相同，与所走路径
无关。也就是说，到达 $B$ 点时的动能等于 $A$ 点的动能与从 $A$ 到 $B$ 的重力功之和；
轨道只改变速度方向，并不改变这个能量关系。

建立竖直向上的 $y$ 轴，设 $A,B$ 两点纵坐标为 $y_A,y_B$，则
\begin{equation}
 h=y_A-y_B.
\end{equation}
代入式~\eqref{ch08:eq:gravity-kinetic}：
\begin{equation}
 \frac12mv_B^2=\frac12mv_A^2+mg(y_A-y_B).
\end{equation}
进一步整理为
\begin{equation}
 \frac12mv_B^2+mgy_B
 =\frac12mv_A^2+mgy_A.
\end{equation}
定义物体在重力场中的重力势能为
\begin{equation}
 E_p(y)=mgy.\label{ch08:eq:gravity-potential}
\end{equation}
上述推理并不依赖 $A,B$ 两点处在某两个特殊高度，因此对运动中任意选取的两个位置都成立。
于是，当物体只受重力以及不做功的约束力时，动能与重力势能之和保持不变。定义机械能
\begin{equation}
 E=E_k+E_p=\frac12mv^2+mgy,
\end{equation}
则整个运动过程中 $E$ 为常量。势能零点可以任意选择；改变零点只会给所有势能同时加上
一个常数，不影响势能差和运动规律。

\begin{origexamplebox}{例 4.10}
质量均为 $m$ 的质点，初速率均为 $v$。轨道固定、无摩擦，不计空气阻力。
高度从初始位置算起，圆周运动均由最低点开始。求以下各情况下首次上升所能达到的最大高度：
\begin{enumerate}[label=\arabic*.,leftmargin=2.2em]
 \item 沿一条高度单调上升且无限高的弯曲光滑槽运动，槽保证质点不脱离；
 \item 以与水平面夹角为 $\theta$ 的方向斜向抛出，$0\le\theta\le\pi/2$；
 \item 套在半径为 $R$ 的光滑圆环上运动；
 \item 在半径为 $R$ 的光滑圆柱体内部运动。
\end{enumerate}
\probdiagFourTen
\end{origexamplebox}

\begin{xesolution}
图~\ref{ch08:fig:probdiagFourTen} 的第三种是套住细圆环的珠子，第四种是内壁单侧接触，不能把二者当作同一约束。
以下高度均从初始位置所在的水平面算起。

\textbf{(1) 无限高的光滑轨道}\quad 物体受到的支持力不做功。到达最高点时速率为零，
\[
 \frac12mv^2=mgH_1,
 \qquad
 \boxed{H_1=\frac{v^2}{2g}}.
\]

\textbf{(2) 斜抛}\quad 重力没有水平分量，所以 $v_x=v\cos\theta$ 保持不变。
最高点只有竖直分速度为零，并非总速度为零。将最高点仍有的水平动能扣除：
\[
 mgH_2=\frac12mv^2-\frac12m(v\cos\theta)^2=\frac12m(v\sin\theta)^2.
\]
因此
\[
 \frac12m(v\sin\theta)^2=mgH_2,
 \qquad
 \boxed{H_2=\frac{v^2\sin^2\theta}{2g}}.
\]

\textbf{(3) 套在圆环上}\quad 圆环可同时提供指向圆心或背离圆心的约束力，质点不会
脱离圆环。这是不同状态下可出现的两种方向，并非同时存在两个支持力。
若初始动能不足以到达环顶，它在中途反向；若能到达环顶，几何上的最大高度
就是 $2R$：
\[
 \boxed{H_3=\min\left(\frac{v^2}{2g},\,2R\right)}.
\]

临界情况 $v^2=4gR$ 下，理想质点无限接近环顶，到达时间趋于无穷；此时 $2R$ 严格说是高度上确界。

\textbf{(4) 在圆柱体内部}\quad 这一情形的约束力只能指向圆心，物体可能脱离内壁。
令质点沿内壁从最低点上升的圆心角为 $\varphi$，并记
\[
 q=\frac{v^2}{gR}.
\]
在接触阶段，能量守恒和径向方程分别为
\[
 u^2=v^2-2gR(1-\cos\varphi),
 \qquad
 N=m\left(\frac{u^2}{R}+g\cos\varphi\right).
\]
代入可得 $N=mg(q-2+3\cos\varphi)$。若 $q\le2$，速度先在下半圆降为零；
若 $q\ge5$，环顶仍有 $N=mg(q-5)\ge0$，能保持接触。
所以仅当 $2<q<5$ 时，物体在上半圆由 $N=0$ 脱离。令
\[
 c=\cos\varphi_{\rm d}=\frac{2-q}{3},
\]
则脱离速度满足 $u_{\rm d}^2=-gRc$。图~\ref{ch08:fig:soldiagFourTen} 的右图中，$D$ 为脱离点，
速度沿切线，竖直分量为 $u_{{\rm d},y}=u_{\rm d}\sin\varphi_{\rm d}$。之后作斜抛，
向上速度降为零前增加的高度为 $u_{{\rm d},y}^2/(2g)$。与脱离点高度相加，
得到
\[
 H_4
 =R(1-c)+\frac{u_{\rm d}^2(1-c^2)}{2g}
 =R\left(1-\frac32c+\frac12c^3\right).
\]
\soldiagFourTen
综上，
\[
 \boxed{
 H_4=
 \begin{cases}
 \dfrac{v^2}{2g},&0\le q\le2,\\[6pt]
 R\left(1-\dfrac32c+\dfrac12c^3\right),
   &2<q<5,\quad c=\dfrac{2-q}{3},\\[6pt]
 2R,&q\ge5.
 \end{cases}}
\]
第一段中质点在到达圆环侧面以前已经反向；中间一段先脱离再作抛体；最后一段始终接触
内壁并越过最高点。在 $q=2,5$ 处相邻表达式分别都给出 $R,2R$，结果连续，每项均为长度。
这里问首次上升；如果再碰到内壁，后续还与碰撞模型有关。只用能量写 $H=v^2/(2g)$ 会漏掉约束失效后的运动。
\end{xesolution}

\Needspace{0.34\textheight}
\begin{origexamplebox}{例 4.11}
一个质量为 $m$ 的质点从固定在水平面上、半径为 $R$ 的半圆柱形光滑表面的顶端由静止
开始下滑。求质点脱离圆柱体时，质点与圆心连线和竖直方向的夹角 $\theta$；并求落地点
与圆心的距离。理想的正顶点静止状态是非稳定平衡；“开始下滑”按微小扰动趋于零的极限理解。
\probdiagFourEleven
\end{origexamplebox}

\begin{xesolution}
图~\ref{ch08:fig:probdiagFourEleven} 中半圆的圆心在地面，顶点距圆心确为 $R$。
设质点在圆柱面上运动到圆心角 $\theta$ 时速率为 $u$。从顶端由静止下滑，机械能守恒：
\[
 \frac12mu^2=mgR(1-\cos\theta),
 \qquad
 u^2=2gR(1-\cos\theta).
\]
取指向圆心为径向正方向。质点在圆柱外表面运动时
\[
 mg\cos\theta-N=\frac{mu^2}{R}.
\]
脱离瞬间 $N=0$，所以
\[
 g\cos\theta=\frac{u^2}{R}.
\]
与能量式联立：
\[
 \cos\theta=2(1-\cos\theta),
 \qquad
 \boxed{\cos\theta=\frac23},
\]
即
\[
 \boxed{\theta=\arccos\frac23}.
\]
此时
\[
 \sin\theta=\frac{\sqrt5}{3},
 \qquad
 u=\sqrt{\frac{2gR}{3}}.
\]
\soldiagFourEleven
图~\ref{ch08:fig:soldiagFourEleven} 以圆心为原点、水平向右为 $x$ 轴、竖直向上为 $y$ 轴。
半径方向为 $(\sin\theta,\cos\theta)$，向右下方的单位切向量为 $(\cos\theta,-\sin\theta)$，二者点积为零。
所以脱离点坐标和速度分量为
\[
 x_0=\frac{\sqrt5}{3}R,\quad
 y_0=\frac23R,\quad
 u_x=\frac23\sqrt{\frac{2gR}{3}},\quad
 u_y=-\frac{\sqrt5}{3}\sqrt{\frac{2gR}{3}}.
\]
脱离后水平速度不变、竖直加速度为 $-g$。落地时 $y=0$，故
\[
 \frac12gt^2-u_yt-y_0=0,\qquad t=\frac{u_y+\sqrt{u_y^2+2gy_0}}{g}.
\]
另一根为负，不能表示脱离后的时刻。代入 $u_y$ 和 $y_0$，得到
\[
 t=\sqrt{\frac{R}{g}}\,
 \frac{\sqrt{46}-\sqrt{10}}{3\sqrt3}.
\]
因此落地点到圆心的水平距离为
\[
 \boxed{
 x=x_0+u_xt
 =\frac{5\sqrt5+4\sqrt{23}}{27}\,R}.
\]
结果约为 $1.12458R$，大于半圆右端的横坐标 $R$，量纲也是长度。
脱离后支持力消失，必须改用抛体运动，不能把圆周运动方程继续使用到落地。
\end{xesolution}

\Needspace{0.33\textheight}
\begin{origexamplebox}{例 4.12}
如图，光滑水平面上质量为 $m$ 的物体以初速度 $v$ 运动，冲向一个半径为 $R$ 的光滑
圆形表面。图~\ref{ch08:fig:probdiagFourTwelve} 的圆弧为固定的右半圆内壁，与水平面在最低点相切。
求使物体脱离圆弧后能够击中圆心 $O$ 的初速度 $v$。
\probdiagFourTwelve
\end{origexamplebox}

\begin{xesolution}
以圆心 $O$ 为原点，令物体从圆弧最低点开始运动，半径由竖直向下转过角
$\varphi$。接触点坐标为
\[
 x=R\sin\varphi,\qquad y=-R\cos\varphi.
\]
设脱离时速率为 $u$。物体在圆弧内侧运动，能量守恒和脱离条件为
\[
 u^2=v^2-2gR(1-\cos\varphi),
\qquad
 N=m\left(\frac{u^2}{R}+g\cos\varphi\right)=0.
\]
因此
\begin{equation}
 u^2=-gR\cos\varphi,\qquad v^2=gR(2-3\cos\varphi).
 \label{ch08:eq:circle-entry}
\end{equation}
脱离点在右上象限，故 $\cos\varphi<0$。切向速度分量为
\[
 u_x=u\cos\varphi,\qquad u_y=u\sin\varphi.
\]
若抛体经过圆心，则存在 $t>0$ 使
\[
 0=R\sin\varphi+u\cos\varphi\,t,
\]
\[
 0=-R\cos\varphi+u\sin\varphi\,t-\frac12gt^2.
\]
\soldiagFourTwelve
图~\ref{ch08:fig:soldiagFourTwelve} 中 $D$ 在右上方、速度指向左上方，抛出后水平分速度保持为负。
令 $c=\cos\varphi<0$、$s=\sin\varphi>0$。由横坐标方程得
\[
 t=-\frac{Rs}{uc}>0.
\]
代入纵坐标方程，除以 $R$，再用 $u^2=-gRc$：
\[
 0=-c-\frac{s^2}{c}-\frac{gRs^2}{2u^2c^2}
   =-c-\frac{s^2}{c}+\frac{s^2}{2c^3}.
\]
乘以 $2c^3$ 并代入 $s^2=1-c^2$，逐项整理：
\[
 0=-2c^4-2s^2c^2+s^2
 =-2c^4-2(1-c^2)c^2+1-c^2=1-3c^2.
\]
因此
\[
 1-3\cos^2\varphi=0.
\]
结合 $\cos\varphi<0$：
\[
 \cos\varphi=-\frac1{\sqrt3}.
\]
代入式~\eqref{ch08:eq:circle-entry}，得到
\[
 \boxed{v=\sqrt{(2+\sqrt3)\,gR}}.
\]
结果中 $v^2/(gR)=2+\sqrt3$ 在 $2$ 与 $5$ 之间，符合在上半圆脱离的条件；$\sqrt{gR}$ 的量纲为速度。
这里“击中圆心”发生在脱离后的抛体阶段；若不先判定脱离条件，不能确定正确的初速度。
\end{xesolution}

\subsection{保守力、势能}

重力做功只与始末高度有关。具有同样性质的力称为保守力。保守力的关键性质可整理为：
\begin{enumerate}[leftmargin=2.2em]
 \item 保守力的大小和方向只由位置决定，不显含时间或速度。近地面重力、固定端理想弹簧
 弹力，以及以后学习的万有引力、静电场中的电场力，都是典型例子。
 \item 保守力做功只与初、末位置有关，与实际路径无关。重力功只取决于高度差，弹簧弹力
 的功只取决于始、末形变量。
 \item 由路径无关性立即得到：物体沿任意闭合路径运动一周，保守力的总功必为零。
 \item 对保守力可以定义只依赖位置的势能。势能不依赖物体的速度、运动时间或到达该位置
 的路径，并且
 \[
 W_{A\to B}=E_p(A)-E_p(B)=-\Delta E_p.
 \]
 势能的形式由相应的保守力决定，势能零点可以自由选取；真正影响运动的是势能差。
 \item 已知保守力随位置的关系并选定势能零点，可以求出各位置的势能；反过来，已知势能
 也可以求出力的大小和方向。空间中势能相等的点构成等势面，保守力与等势面垂直，并指向
 势能下降最快的方向。
 \item 若物体只受保守力和不做功的约束力作用，便可把机械能定义为动能与各项势能之和，
 且运动过程中机械能保持不变；这个常量可由运动中任意一个位置的动能和势能确定。
\end{enumerate}
重力势能已由式~\eqref{ch08:eq:gravity-potential} 给出。劲度系数为 $k$ 的理想弹簧，形变量为 $x$ 时，弹性势能为
\begin{equation}
 E_p(x)=\frac12kx^2.
\end{equation}
对于万有引力，若把无穷远处选作势能零点，两质量 $M,m$ 相距 $r$ 时的势能为
\begin{equation}
 E_p(r)=-\frac{GMm}{r}.
\end{equation}
势能的绝对值取决于零点选择，真正决定运动的是势能差。一维微小位移满足
$F_x\dd x=-\dd E_p$，故 $F_x=-\dd E_p/\dd x$：曲线斜率为正时力向负方向，斜率为负时力向正方向。
三维情况下，势能一般是 $E_p(x,y,z)$，全微分为
\[
 \dd E_p=\frac{\partial E_p}{\partial x}\dd x+
 \frac{\partial E_p}{\partial y}\dd y+\frac{\partial E_p}{\partial z}\dd z.
\]
例如 $\partial E_p/\partial x$ 表示保持 $y,z$ 不变时，势能随 $x$ 的变化率。
与 $\vec F\cdot\dd\vec r=-\dd E_p$ 逐项比较，三个分量分别是对应偏导数的负值。
把三个偏导数组成矢量，称为势能的梯度，记为
\[
 \nabla E_p=\left(\frac{\partial E_p}{\partial x},\frac{\partial E_p}{\partial y},\frac{\partial E_p}{\partial z}\right),
 \qquad \vec F=-\nabla E_p.
\]
在梯度非零的地方，梯度指向势能增加最快的方向；因此保守力与等势面垂直，指向势能降低最快的方向。
若物体只受保守力作用，则动能和相应势能之和守恒，即
\[
 E=E_k+E_p=\text{常量}.
\]

\Needspace{22\baselineskip}
\begin{origexamplebox}{例 4.13}
判断下列各力是否为保守力；若是，分析其势能函数的性质。各小题独立，系数取正值；
$a$ 在第(1)问是力随时间增加的系数，不是加速度，各题 $\alpha$ 的单位由所在公式决定。
\begin{enumerate}[label=\arabic*.,leftmargin=2.2em]
 \item 正比于时间的拉力 $f=a t$；
 \item 正比于速度、方向与速度相反的阻力 $f=-\alpha v$；
 \item 正比于坐标三次方的力 $f=-kx^3$；
 \item 随距离变化的吸引力
 \[
 f(r)=-k\frac{e^{-\alpha r}}{r^2}.
 \]
\end{enumerate}
\end{origexamplebox}

\begin{xesolution}
\textbf{(1)} $f=at$ 显含时间。同一位置在不同时刻受到的力不同，不能由单值的位置势能
确定，因而不是通常意义下的保守力。

\textbf{(2)} $f=-\alpha v$ 依赖速度。其瞬时功率为
\[
 P=fv=-\alpha v^2\le0,
\]
机械能随运动永久耗散，故不是保守力。

\textbf{(3)} $f=-kx^3$ 只依赖位置，在一维中为保守力。由
\[
 -\frac{\dd E_p}{\dd x}=-kx^3
\]
得到
\[
 \boxed{E_p(x)=\frac{k}{4}x^4+C}.
\]
若取 $E_p(0)=0$，则 $C=0$。势能在 $x=0$ 取得稳定极小值，关于原点对称。
但最低点的二阶导数为零，小振动不能直接当作具有非零固有频率的简谐振动。

\textbf{(4)} 这是只依赖距离的中心力，在 $r>0$ 的区域内为保守力。取
$E_p(\infty)=0$，可写成
\[
 \boxed{
 E_p(r)=-\int_r^\infty
 k\frac{e^{-\alpha s}}{s^2}\,\dd s}.
\]
上述定积分已经完整定义了势能，不必依赖特殊函数。若进一步引入指数积分函数
\[
 \operatorname{Ei}(-z)=-\int_z^\infty\frac{e^{-t}}{t}\,\dd t,\qquad z>0,
\]
注意 $\dd(e^{-\alpha s}/s)/\dd s=-\alpha e^{-\alpha s}/s-e^{-\alpha s}/s^2$，因此分部积分可得
\[
 \int_r^\infty\frac{e^{-\alpha s}}{s^2}\dd s
 =\frac{e^{-\alpha r}}r-\alpha\int_r^\infty\frac{e^{-\alpha s}}s\dd s
 =\frac{e^{-\alpha r}}r+\alpha\operatorname{Ei}(-\alpha r).
\]
于是也可表示为
\[
 E_p(r)
 =-k\left[
 \frac{e^{-\alpha r}}{r}
 +\alpha\,\operatorname{Ei}(-\alpha r)
 \right].
\]
\soldiagFourThirteen
图~\ref{ch08:fig:soldiagFourThirteen} 的右图用这个定积分作图，而不是 $-ke^{-\alpha r}/r$；后者求负导数会多出一项，不能代表题给的力。
由微积分基本定理，$E_p'(r)=ke^{-\alpha r}/r^2>0$，负导数恰为题给的引力。
它随 $r$ 单调增加：$r\to\infty$ 时趋于 $0$，$r\to0^+$ 时趋于 $-\infty$，对应始终
指向内侧的吸引力。前三项中，只有第 (3) 项具有普通的位置势能；第 (4) 项虽形式复杂，
但是否保守只取决于它是位置的单值中心力，而不取决于积分能否用初等函数表示。
\end{xesolution}

\Needspace{0.32\textheight}
\begin{origexamplebox}{例 4.14}
劲度系数为 $k$、自然长度为 $l_0$ 的弹簧一端固定在给定点 $A$，另一端连接质量为
$m$ 的物体。前两问限定物体沿弹簧轴线作一维运动，弹簧可拉伸也可压缩，忽略碰壁等几何限制。
求以下各情况下弹簧的最大伸长量：
\begin{enumerate}[label=\arabic*.,leftmargin=2.2em]
 \item 物体在与 $A$ 同一水平面的光滑平面上运动。初始时弹簧为原长，物体初速度为
 $v_1$；
 \item 仍在上述光滑水平面内运动。初始时弹簧压缩量为 $\delta l$，初速度为 $v_2$，
 方向指向平衡位置。若初速度方向背离平衡位置，结果如何？
 \item 物体在通过 $A$ 的竖直线上运动。初始时位于原长位置以下 $h$ 处，速度 $v_3$
 竖直向上。若初速度改为竖直向下，结果如何？
\end{enumerate}
\probdiagFourFourteen
\end{origexamplebox}

\begin{xesolution}
图~\ref{ch08:fig:probdiagFourFourteen} 将水平轴线和竖直轴线分开表示，固定点都记为 $A$。
以下用 $x$ 表示弹簧相对原长的形变量，规定伸长为正；“最大伸长量”就是运动中最大的
$x$。

\textbf{(1)} 水平面上没有重力势能变化。初态 $x=0$，在最大伸长处速度为零：
\[
\frac12mv_1^2=\frac12kx_{\max}^2,
\qquad
\boxed{x_{\max}=v_1\sqrt{\frac{m}{k}}}.
\]

\textbf{(2)} 初态为 $x=-\delta l$。在正向最大伸长处速度为零，故
\[
\frac12mv_2^2+\frac12k(\delta l)^2
=\frac12kx_{\max}^2,
\qquad
\boxed{
 x_{\max}=\sqrt{(\delta l)^2+\frac{m}{k}v_2^2}}.
\]
初速度若背离平衡位置，物体先到压缩侧的转折点，再向另一侧运动；由于总机械能不变，
最终的最大伸长量仍相同，只有到达各转折点的先后次序改变。

\textbf{(3)} 取竖直向下为 $x$ 正方向，忽略与原长 $l_0$ 有关的势能常数：
\[
E_p(x)=\frac12kx^2-mgx
=\frac12k\left(x-\frac{mg}{k}\right)^2+\text{常量}.
\]
竖直运动的平衡形变量和相对该位置的振幅分别为
\[
x_{\rm e}=\frac{mg}{k},
\qquad
\mathcal A=\sqrt{\left(h-\frac{mg}{k}\right)^2+\frac{m}{k}v_3^2}.
\]
\soldiagFourFourteen
这里 $\mathcal A$ 表示振幅，区别于固定点 $A$。它来自初态与转折点的能量比较：
\[
 \frac12mv_3^2+\frac12k(h-x_{\rm e})^2=\frac12k\mathcal A^2.
\]
图~\ref{ch08:fig:soldiagFourFourteen} 中势阱关于 $x_{\rm e}$ 对称，两个转折点为 $x_{\rm e}\pm\mathcal A$。
最大伸长位置为 $x_{\rm e}+\mathcal A$，所以
\[
\boxed{
 x_{\max}=\frac{mg}{k}
 +\sqrt{\left(h-\frac{mg}{k}\right)^2+\frac{m}{k}v_3^2}}.
\]
把初速度由向上改为向下，只改变速度符号，不改变机械能和振幅，故最大伸长量仍不变。
取 $v_3=0,h=x_{\rm e}$ 时振幅为零，物体保持平衡，结果给出 $x_{\max}=x_{\rm e}$。
若前两问不限制为沿轴线运动，非径向初速度会产生角动量；最大伸长处径向速度虽为零，切向速度却不一定为零。
此时需另用角动量守恒和包含切向动能的能量式，不能直接使用本题的一维答案。
\end{xesolution}

\subsection{势能方法的适用范围与稳定性}

势能函数也常记为 $U$。近地面重力势能 $mgy$、万有引力势能 $-GMm/r$、弹簧势能 $kx^2/2$
来自不同的力随位置的关系，不能跨越适用范围混用。尤其当高度与地球半径可比时，$g$ 随位置变化，$mgy$ 不再准确。

先选出已用势能 $U$ 表示的、与时间无关的保守相互作用；把所有尚未计入 $U$ 的力的功合记为 $W_{\rm other}$。
动能定理给出 $\Delta K=-\Delta U+W_{\rm other}$，所以
\begin{equation}
 \Delta(K+U)=W_{\rm other}.
 \label{ch08:eq:other-work}
\end{equation}
当 $W_{\rm other}=0$，机械能守恒。若还要把右边分为外力功和非保守内力功，两类必须互不重叠，
不能把同一个外部摩擦力计算两次。也不能既写重力功，又把同一重力势能变化列入左边。
在固定外场中可以把外场作用写入势能；例如研究近地面质点时使用 $mgy$，不必每次显式计算地球的运动。
但若场源也运动或相互作用显含时间，就须重新检查这个简化是否成立。

一维平衡点 $x_0$ 满足 $U'(x_0)=0$，即合力为零。若 $U''(x_0)>0$，势能在该点取局部极小，
微小偏离后出现指向平衡点的回复力，是稳定平衡；若 $U''(x_0)<0$，回复方向相反，平衡不稳定。
在二阶导数为正的平衡点附近展开：
\[
 U(x)\simeq U(x_0)+\frac12U''(x_0)(x-x_0)^2,
 \qquad F_x\simeq-U''(x_0)(x-x_0).
\]
与 $m\ddot x=-k_{\rm eff}(x-x_0)$ 比较可知，$k_{\rm eff}=U''(x_0)$，故小振动角频率
$\omega=\sqrt{U''(x_0)/m}$。如果 $U''(x_0)=0$，这个二阶判据不能定性，必须继续查看更高阶项；
例如 $U=kx^4/4$ 仍是稳定势阱，却不是上述线性回复模型。

给定一维总机械能 $E$，允许位置满足 $E\ge U(x)$；转折点满足 $E=U(x)$。在允许区间内
\[
 |v(x)|=\sqrt{\frac{2[E-U(x)]}{m}}.
\]
这只能先给出速率，运动方向由所处的往返分支确定。若需时间，在方向不变的一段上取
$v(x)=\pm\sqrt{2[E-U(x)]/m}$，再积分 $t_B-t_A=\int_{x_A}^{x_B}\dd x/v(x)$。
因此能量曲线既能判断可达区域，也能与运动方程配合求运动时间。

